\documentclass[../../../include/open-logic-section]{subfiles}

\begin{document}

\olfileid{cmp}{rec}{pre}
\olsection{原始递归}

% Part: computability
% Chapter: recursive-functions
% Section: primitive-recursion
自然数的一个典型特征是：每个自然数都可以从~$0$ 出发，通过有限次应用后继运算~$+1$ 而得到——任何自然数要么是~$0$，要么是~\dots{}~$0$ 的后继。利用这一事实定义函数~$h\colon\Nat \to \Nat$ 的一种方式是：(a)~指定 $h$ 在自变元~$0$ 处的值，并且 (b)~说明如何根据 $h(x)$ 的值去计算 $h(x+1)$ 的值。因为 (a) 直接告诉我们 $h(0)$ 是什么，所以 $h$ 在~$0$ 处有定义。现在，利用 (b) 对 $x=0$ 给出的指示，我们可以从~$h(0)$ 计算出 $h(1) = h(0+1)$。再对 $x=1$ 应用同样的指示，我们从~$h(1)$ 计算出 $h(2) = h(1+1)$，依此类推。对于每个自然数~$x$，我们最终都会到达从 $h(x+1)$ 定义 $h(x)$ 的那一步，因此 $h(x)$ 对所有的 $x \in \Nat$ 都有定义。

例如，假设我们通过以下两个方程来指定 $h\colon \Nat \to \Nat$：
\begin{align*}
h(0) & =  1\\
h(x+1) & =  2 \cdot h(x)
\end{align*}
如果我们已经知道如何做乘法，那么这两个方程就给出了上面 (a) 和~(b) 所要求的信息。连续应用第二个方程，我们得到：
\begin{align*}
  h(1) & = 2\cdot h(0) = 2,\\
  h(2) & = 2\cdot h(1) = 2\cdot 2,\\
  h(3) & = 2 \cdot h(2) = 2\cdot 2 \cdot 2,\\
  & \vdots
\end{align*}
可以看出，我们所指定的函数~$h$ 就是 $h(x) = 2^x$。

自然数的这一典型特征保证了只有一个函数~$h$ 满足这两条准则。像这样的一对方程称为函数~$h$ 的一个\emph{原始递归定义}。之所以这么称呼，是因为我们“递归地”定义~$h$，即定义式（具体是第二个方程）的右边涉及 $h$ 自身。称之为“原始”的，是因为在定义 $h(x+1)$ 时，我们只使用了值~$h(x)$，即紧邻的前一个值。这是在~$\Nat$ 上递归定义函数的最简单方式。

我们甚至可以用原始递归来定义更基本的函数，比如加法和乘法。不过在这些情况下，所讨论的函数是 $2$~元的。我们固定其中一个自变元的位置，用另一个位置进行递归。例如，要定义 $\Add(x, y)$，我们可以固定~$x$，先定义 $y=0$ 时的值，然后根据~$y$ 来定义 $y+1$ 时的值。由于 $x$ 是固定的，它会出现在定义方程的左边和右边。
\begin{align*}
\Add(x,0) & =  x\\
\Add(x,y+1) & =  \Add(x,y)+1
\end{align*}
这些方程指定了 $\Add$ 对所有 $x$ \emph{和}~$y$ 的值。例如，要计算 $\Add(2,3)$，我们就对 $x = 2$ 应用定义方程，先用第一个方程得到 $\Add(2,0) = 2$，然后用第二个方程依次求出 $\Add(2,1) = 2 + 1 = 3$、$\Add(2, 2) = 3 + 1 = 4$、$\Add(2,
3) = 4 + 1 = 5$。

在 $\Add$ 的定义中，我们在第二个方程的右侧用了 $+$，但只是为了加~$1$。换句话说，我们使用了后继函数 $\Succ(z) = z+1$，并将其应用于前一个值 $\Add(x,y)$ 来定义 $\Add(x,y+1)$。因此，我们可以认为该递归定义是通过一个单一的函数（我们将它应用于前一个值）来给出的。然而，允许该函数不仅依赖前一个值，还依赖 $x$ 和~$y$，这并无妨害——有时甚至是必要的。考虑：
\begin{align*}
  \Mult(x,0) & =  0 \\
  \Mult(x,y+1) & =  \Add(\Mult(x,y),x)
\end{align*}
这是函数 $\Mult$ 的一个原始递归定义，它是通过将函数 $\Add$ 同时应用于前一个值 $\Mult(x,y)$ 和第一个自变元~$x$ 而给出的。它同样为所有自变元 $x$ 和~$y$ 定义了函数~$\Mult(x,y)$。例如，$\Mult(2,3)$ 是通过依次计算 $\Mult(2,0)$、$\Mult(2,1)$、$\Mult(2,2)$ 和~$\Mult(2,3)$ 来确定的：
\begin{align*}
  \Mult(2,0) & = 0\\
  \Mult(2,1) & = \Mult(2,0+1) =
  \Add(\Mult(2,0), 2) = \Add(0, 2) = 2\\
  \Mult(2,2) & = \Mult(2,1+1) =
  \Add(\Mult(2,1), 2) = \Add(2, 2) = 4\\
  \Mult(2,3) & = \Mult(2,2+1) =
  \Add(\Mult(2,2), 2) = \Add(4, 2) = 6
\end{align*}

因此，一般模式如下：要给出一个函数~$h(x_0, \dots, x_{k-1}, y)$ 的原始递归定义，我们提供两个方程。第一个方程不依赖于~$h$，定义了 $h(x_0, \dots, x_{k-1}, 0)$ 的值。第二个方程则根据 $h(x_0, \dots, x_{k-1}, y)$、其他自变元 $x_0$、\dots、$x_{k-1}$ 以及~$y$ 来定义 $h(x_0,
\dots, x_{k-1}, y+1)$ 的值。在第二个方程中，只能使用~$h$ 的紧邻的前一个值。如果我们把这两个方程右边给出的运算本身看作函数 $f$ 和~$g$，那么通过原始递归定义新函数~$h$ 的一般模式就是：
\begin{align*}
  h(x_0, \dots, x_{k-1}, 0) & = f(x_0, \dots, x_{k-1})\\
  h(x_0, \dots, x_{k-1}, y+1) & =
  g(x_0, \dots, x_{k-1}, y, h(x_0, \dots, x_{k-1}, y))
\end{align*}
在 $\Add$ 的情形，我们有 $k=1$，$f(x_0) = x_0$（恒等函数），以及 $g(x_0, y, z) = z + 1$（这个 $3$ 元函数返回其第三个自变元的后继）：
\begin{align*}
  \Add(x_0, 0) & = f(x_0) = x_0\\
  \Add(x_0, y+1) & = g(x_0, y, \Add(x_0, y)) =
  \Succ(\Add(x_0, y))
\end{align*}
在 $\Mult$ 的情形，我们有 $f(x_0) = 0$（总是返回~$0$ 的常值函数）和 $g(x_0, y, z) = \Add(z,x_0)$（这个 $3$ 元函数返回其最后一个自变元与第一个自变元之和）：
\begin{align*}
  \Mult(x_0,0) & =  f(x_0) = 0 \\
  \Mult(x_0,y+1) & =  g(x_0, y, \Mult(x_0,y)) =
  \Add(\Mult(x_0,y), x_0)
\end{align*}

\end{document}